\section{Methods}
\subsection{Inference: a binomial model of the odds of robbery}
For incident $i$ with outcome $y_i=1$ if robbery, we fit a logistic regression with hour of day (reference: noon), month (reference: January), weekday (reference: Monday), an indicator that the sun was below civil twilight at the middle of the local hour, and an indicator for the post-2019 period (occurrence year 2020 or later). Darkness comes from NOAA solar equations for Toronto's coordinates, so unlike a fixed clock cut-off it follows the seasons; with hour and month as factors it is identified only by twilight hours that are dark on some dates and light on others. Hour and month are factors rather than linear numbers, and the factor coding fits better than a cyclic-harmonic alternative (AIC higher by \resG{fit.harmonic.delta_aic} for the harmonic model). The model has \resG{fit.parameters} parameters and uses \resG{n} rows. Incidents within a neighbourhood are not independent, so we report odds ratios with \pct{} intervals from cluster-robust standard errors \citep{cameron2015cluster} over \resG{log.neighbourhoods} neighbourhoods. As a discrimination check, a seeded five-fold cross-validated AUC is computed. The ``post-2019'' term absorbs every change after 2019 (reporting, policing, the mix of other crimes) and is not an effect of any single cause.

\subsection{Prediction: an ablation ladder}
Five models, called rungs, are fitted on identical rows, each adding one ingredient. \textbf{L1} uses time only: cyclic hour and month terms, weekday and darkness, in a regularised logistic regression. \textbf{L2} adds place: premises type, location type and a spline of the coordinates. \textbf{L3} adds the nine census features. \textbf{L4} replaces the logistic model with a random forest \citep{breiman2001rf} and \textbf{L5} with gradient-boosted trees (XGBoost, \citealt{chen2016xgboost}). A constant-score baseline, L0, has ROC-AUC one half and a PR-AUC equal to the share of robberies. Because premises and location type nearly define some crime types (a break and enter happens in a building), every rung from L2 up is also fitted without them. Hyperparameters are chosen from small grids by PR-AUC in a three-fold neighbourhood-grouped cross-validation inside each training part.

\subsection{Evaluation}
We use two designs, each producing out-of-sample predictions for every row. \textbf{Design A} is five-fold cross-validation grouped by neighbourhood, so each neighbourhood's census values and nearby coordinates never appear in both training and test parts \citep{roberts2017cv}; it measures performance in new neighbourhoods. \textbf{Design B} trains on incidents through 2021 and tests on 2022 and 2023, the period when the robbery share was lower; it measures performance in later years. Models are scored on the pooled predictions by ROC-AUC, PR-AUC and F1, plus precision and recall, at the threshold that maximises F1 on inner out-of-fold predictions; the test rows never choose the threshold. PR-AUC depends on prevalence \citep{saito2015pr}, which differs between designs (\resA{prevalence} and \resB{prevalence}), so values are compared only within a design. Intervals are percentile intervals from a cluster bootstrap over neighbourhoods \citep{efron1993bootstrap} with \resA{bootstrap_replicates} replicates and a fixed seed; differences between rungs use the same resamples (paired). The importance of a feature block is the PR-AUC lost on the held-out rows when the block's columns are shuffled jointly, for the L5 model under Design A.

\subsection{Corrections to the course versions}
The two earlier analyses were re-run before being changed. Relative to them: the course regression's final summary table was accidentally a Gaussian linear model, so its coefficients were not odds ratios; hour, weekday and month entered as linear numbers, darkness was a fixed clock rule and standard errors ignored clustering; the reported AUC was apparent rather than cross-validated; the classifier used one random row split, so neighbourhoods appeared in both training and test parts, a fixed threshold of one half, and one year of data; and the written XGBoost tuning grid was never passed to the training call, so a smaller default grid was searched. The versions here fix each of these, seed every random step and use one set of rows for every model.
